\section{DICOM và khả năng liên thông trong ảnh y khoa}
\label{sec:dicom}

Bất kỳ nền tảng nào muốn vận hành bên trong bệnh viện đều phải nói được ngôn ngữ
của bệnh viện. Đối với chẩn đoán hình ảnh, ngôn ngữ đó là DICOM, và các đặc tính
của nó ràng buộc kiến trúc hệ thống đủ nhiều để cần được trình bày như một nội
dung riêng thay vì chỉ là một chi tiết cài đặt.

\subsection{Chuẩn DICOM}

DICOM (Digital Imaging and Communications in Medicine, ISO 12052) là chuẩn quốc
tế quy định định dạng, lưu trữ và trao đổi ảnh y khoa \cite{dicom2024,pianykh2012dicom}. Chuẩn này bắt nguồn từ nỗ
lực của ACR--NEMA năm 1985 nhằm chấm dứt tình trạng phụ thuộc nhà cung cấp
(vendor lock-in) giữa các hãng sản xuất thiết bị chẩn đoán hình ảnh, và đã trở
nên phổ quát trên thực tế kể từ phiên bản 3.0 (1993). Điểm quan trọng là DICOM
không chỉ đơn thuần là một định dạng tệp. Nó đồng thời đặc tả:

\begin{itemize}
    \item một \textbf{mô hình thông tin} (information model), tổ chức dữ liệu
          theo phân cấp bốn mức Bệnh nhân $\rightarrow$ Ca chụp (Study)
          $\rightarrow$ Chuỗi ảnh (Series) $\rightarrow$ Đối tượng ảnh
          (Instance);
    \item một \textbf{định dạng tệp} trong đó dữ liệu điểm ảnh và siêu dữ liệu
          không thể tách rời, được mang trong cùng một đối tượng; và
    \item một họ \textbf{giao thức mạng} dùng để truy vấn, truy xuất và truyền
          tải các đối tượng đó.
\end{itemize}

Tính không thể tách rời giữa điểm ảnh và siêu dữ liệu là điểm xuất phát cho mọi
vấn đề tiếp theo. Một đối tượng DICOM là một chuỗi các \emph{phần tử dữ liệu}
(data element), mỗi phần tử được định danh bởi một \emph{thẻ} (tag), là cặp số
16 bit viết dưới dạng \texttt{(group,element)}. Bảng~\ref{tab:tags} liệt kê các
phần tử mà đề tài này phụ thuộc vào.

\begin{table}[H]
\centering
\small
\begin{tabularx}{\textwidth}{@{}l l X@{}}
\toprule
\textbf{Thẻ} & \textbf{Tên} & \textbf{Ý nghĩa} \\
\midrule
\texttt{(0010,0010)} & PatientName & Định danh trực tiếp; phải được loại bỏ trước khi dữ liệu rời khỏi bệnh viện \\
\texttt{(0010,0020)} & PatientID & Mã hồ sơ bệnh án của bệnh viện; định danh trực tiếp \\
\texttt{(0010,0030)} & PatientBirthDate & Bán định danh; có nguy cơ tái định danh khi kết hợp \\
\texttt{(0020,000D)} & StudyInstanceUID & Định danh duy nhất toàn cục của ca chụp; khóa chính của một ca chụp \\
\texttt{(0020,000E)} & SeriesInstanceUID & Định danh một chuỗi ảnh trong một ca chụp \\
\texttt{(0008,0018)} & SOPInstanceUID & Định danh một đối tượng ảnh đơn lẻ \\
\texttt{(0008,0050)} & AccessionNumber & Mã chỉ định do HIS/RIS cấp; khoá liên kết ca chụp với chỉ định và kết quả đọc (Mục~\ref{sec:his}) \\
\texttt{(0008,0060)} & Modality & \texttt{CT}, \texttt{MR}, \ldots; dùng để định tuyến \\
\texttt{(0028,1052)} & RescaleIntercept & Chuyển giá trị lưu trữ sang đơn vị Hounsfield \\
\texttt{(0028,1053)} & RescaleSlope & Chuyển giá trị lưu trữ sang đơn vị Hounsfield \\
\texttt{(0028,1050)} & WindowCenter & Cửa sổ hiển thị mặc định \\
\texttt{(0028,1051)} & WindowWidth & Cửa sổ hiển thị mặc định \\
\bottomrule
\end{tabularx}
\caption{Các phần tử dữ liệu DICOM mà nền tảng phụ thuộc vào.}
\label{tab:tags}
\end{table}

Giá trị điểm ảnh được lưu trữ bản thân chúng không phải là đơn vị Hounsfield;
phép chuyển đổi là

\begin{equation}
    \mathrm{HU} = \mathtt{RescaleSlope} \times p + \mathtt{RescaleIntercept},
    \label{eq:hu}
\end{equation}

trong đó $p$ là giá trị được lưu trữ. Áp dụng Công thức~\ref{eq:hu} trước mọi
thao tác cửa sổ là điều kiện để hình ảnh hiển thị nhất quán giữa các máy chụp
khác nhau và để các phép đo đậm độ có ý nghĩa so sánh được, bởi các nhà sản
xuất khác nhau dùng cách mã hóa giá trị lưu trữ khác nhau cho cùng một mức suy
giảm vật lý. Bỏ qua bước này thì cùng một mô, khi chụp trên hai thiết bị khác
hãng, sẽ hiển thị khác nhau trên trình xem.

\subsection{Các dịch vụ mạng}

Các thao tác mạng của DICOM được xây dựng trên các \emph{lớp dịch vụ} (service
class), trong đó một bên đóng vai trò Service Class User (SCU, bên yêu cầu) và
bên kia đóng vai trò Service Class Provider (SCP, bên đáp ứng). Các thao tác
liên quan ở đây gồm:

\begin{description}[leftmargin=2.2cm,style=nextline]
    \item[C-STORE] truyền một đối tượng ảnh từ nút này sang nút khác. Đây là cách
          máy chụp CT gửi một ca chụp đã hoàn tất tới PACS, và là điểm tự nhiên
          để nền tảng chặn bắt các ca chụp mới.
    \item[C-FIND] truy vấn các ca chụp thỏa mãn tiêu chí cho trước.
    \item[C-MOVE] yêu cầu một nút gửi các đối tượng ảnh phù hợp tới một bên thứ ba.
    \item[C-ECHO] kiểm tra kết nối, tương đương lệnh ping trong DICOM.
\end{description}

Các thao tác này chạy trên một liên kết TCP tại cổng thường được đánh số 104 hoặc
11112, và định danh các bên tham gia bằng \emph{Application Entity Title} (AE
Title), một tên ký hiệu ngắn. Việc thiết lập liên kết đòi hỏi mỗi bên phải nhận
biết AE Title của bên kia, qua đó cung cấp một cơ chế kiểm soát truy cập sơ khai:
một nút sẽ từ chối liên kết với một đối tác lạ. Đây là cấu hình chứ không phải
xác thực theo nghĩa mật mã học, và giao thức DICOM cổ điển mặc định truyền dữ
liệu ở dạng rõ. Do đó, mọi triển khai đi qua mạng không tin cậy đều phải bọc
DICOM trong TLS hoặc VPN, một yêu cầu sẽ xuất hiện trở lại ở
Mục~\ref{sec:architecture}.

\subsection{DICOMweb}

DICOMweb là cách diễn đạt lại các dịch vụ trên theo kiểu RESTful trên nền HTTP,
gồm WADO-RS (truy xuất), STOW-RS (lưu trữ) và QIDO-RS (truy vấn) \cite{dicomweb}. Nó quan trọng
với đề tài này vì một lý do thực tiễn: các trình xem chạy trên trình duyệt không
thể mở liên kết mạng DICOM thuần, nên một giao diện lâm sàng dạng web đòi hỏi
DICOMweb ở phía máy chủ. DICOMweb cũng kết hợp tự nhiên với hạ tầng HTTP thông
thường, bao gồm TLS, xác thực bằng bearer-token và reverse proxy, điều mà giao
thức cổ điển không làm được.

Bản mẫu của nền tảng phản ánh trực tiếp điều này. Orthanc, một PACS nhẹ mã nguồn
mở \cite{jodogne2018orthanc}, được cấu hình để đồng thời mở cổng DICOM cổ điển (4242) nhằm tiếp nhận ca
chụp từ các thiết bị và một giao diện DICOMweb (8042) cho trình xem, với reverse
proxy của trình xem chuyển tiếp \texttt{/dicom-web} và \texttt{/wado} tới đó.
Cách bố trí song giao thức này là điều cho phép một giao diện web hiện đại cùng
tồn tại với thiết bị chẩn đoán hình ảnh chỉ nói được giao thức năm 1993.

\subsection{Cú pháp truyền và nén ảnh}
\label{sec:transfersyntax}

Mỗi đối tượng DICOM được mã hoá theo một \emph{cú pháp truyền} (Transfer Syntax),
định danh bằng một UID, quy định thứ tự byte, cách biểu diễn kiểu dữ liệu và,
quan trọng nhất, cách nén dữ liệu điểm ảnh \cite{dicomps35}. Hai bên của một liên
kết DICOM thương lượng cú pháp truyền khi thiết lập liên kết; một nút không hỗ trợ
cú pháp của đối tác sẽ yêu cầu chuyển mã hoặc từ chối đối tượng. Bảng~\ref{tab:ts}
liệt kê các cú pháp liên quan tới đề tài.

\begin{table}[H]
\centering
\small
\begin{tabularx}{\textwidth}{@{}l l X@{}}
\toprule
\textbf{Cú pháp truyền} & \textbf{UID} & \textbf{Tính chất} \\
\midrule
Implicit VR Little Endian & \texttt{1.2.840.10008.1.2} & Mặc định bắt buộc phải hỗ trợ; không nén \\
Explicit VR Little Endian & \texttt{1.2.840.10008.1.2.1} & Không nén; phổ biến nhất \\
JPEG Lossless (SV1) & \texttt{1.2.840.10008.1.2.4.70} & Không mất mát; được thiết bị đời cũ hỗ trợ rộng rãi \\
JPEG-LS Lossless & \texttt{1.2.840.10008.1.2.4.80} & Không mất mát; nhanh, tỉ lệ nén tốt \\
JPEG 2000 (Lossless Only) & \texttt{1.2.840.10008.1.2.4.90} & Không mất mát; mở rộng được theo độ phân giải nhưng tốn tính toán khi giải mã \\
HTJ2K (Lossless Only) & \texttt{1.2.840.10008.1.2.4.201} & Bổ sung năm 2023 (Supplement 235); thông lượng mã hoá và giải mã cao \\
HTJ2K (Lossless RPCL) & \texttt{1.2.840.10008.1.2.4.202} & Như trên, sắp thứ tự luồng bit tối ưu cho hiển thị lũy tiến theo độ phân giải \\
HTJ2K & \texttt{1.2.840.10008.1.2.4.203} & Cho phép cả nén có mất mát lẫn không mất mát \\
\bottomrule
\end{tabularx}
\caption{Các cú pháp truyền DICOM liên quan tới nền tảng.}
\label{tab:ts}
\end{table}

Từ Mục~\ref{sec:windowing}, yêu cầu đối với nền tảng là rõ ràng: bản lưu trữ phục
vụ chẩn đoán chỉ được dùng cú pháp không mất mát. Nén có mất mát làm thay đổi giá
trị điểm ảnh, và vì phép cửa sổ khuếch đại các chênh lệch nhỏ về HU thành chênh
lệch lớn về cường độ hiển thị (Hình~\ref{fig:windowing}), sai số nén vốn khó thấy
trên ảnh thông thường có thể trở nên đáng kể khi đọc trong cửa sổ hẹp. Nén có mất
mát vì vậy chỉ được dùng cho ảnh thu nhỏ và ảnh xem trước. Với ảnh CT, nén không
mất mát thường chỉ giảm dung lượng vài lần, nhưng như vậy đã đủ thay đổi đáng kể
thời gian truyền qua đường truyền chậm (Hình~\ref{fig:transfer}) và dung lượng
lưu trữ dài hạn (Hình~\ref{fig:storage}).

Trong nhóm không mất mát, HTJ2K \cite{dicomsup235} đáng được chú ý vì hai lý do.
Thứ nhất, nó giữ các tính năng của JPEG 2000, như khả năng mở rộng theo độ phân
giải, nhưng thay bộ mã hoá khối bằng một thuật toán có thông lượng cao hơn đáng
kể, loại bỏ nhược điểm chính khiến JPEG 2000 ít được dùng trong hiển thị thời gian
thực. Thứ hai, với biến thể RPCL, trình xem có thể hiển thị một phiên bản độ phân
giải thấp của lát cắt ngay khi nhận được phần đầu của luồng bit và làm nét dần khi
phần còn lại đến. Cornerstone3D, lớp hiển thị của OHIF từ phiên bản 3, hỗ trợ kiểu
tải lũy tiến này, và dịch vụ AWS HealthImaging cũng lưu dữ liệu điểm ảnh dưới dạng
HTJ2K \cite{awshealthimaging}. Với đường truyền chậm ở bối cảnh C, đây là cơ chế
cho phép bác sĩ tuyến trên bắt đầu đọc trước khi toàn bộ ca chụp được truyền xong;
mức lợi ích thực tế cần được đo trong Giai đoạn~2.

\subsection{Tích hợp với hệ thống thông tin bệnh viện}
\label{sec:his}

Ảnh chụp chỉ có ý nghĩa lâm sàng khi gắn với đúng bệnh nhân, đúng chỉ định và đúng
kết quả đọc. Những thông tin này nằm trong hệ thống thông tin bệnh viện (HIS), hệ
thống thông tin chẩn đoán hình ảnh (RIS) và hồ sơ bệnh án điện tử (EMR), không nằm
trong PACS. Một nền tảng quản lý ảnh vì vậy phải trao đổi với các hệ thống đó qua
các cơ chế chuẩn sau.

\begin{description}[leftmargin=2.6cm,style=nextline]
\item[Modality Worklist] Thiết bị chụp truy vấn danh sách ca đã được chỉ định
      (C-FIND trên worklist) để lấy thông tin bệnh nhân, mã chỉ định
      (\texttt{AccessionNumber}) và thủ thuật cần chụp, thay vì kỹ thuật viên
      nhập tay. Đây là biện pháp hiệu quả nhất để tránh sai lệch định danh ngay
      từ nguồn. Modality Performed Procedure Step (MPPS) báo ngược trạng thái
      đã chụp.
\item[HL7 v2] Chuẩn nhắn tin vẫn phổ biến nhất giữa HIS và RIS: bản tin ADT cho
      thông tin hành chính của bệnh nhân, ORM/OMI cho chỉ định, ORU cho kết
      quả.
\item[HL7 FHIR] Chuẩn REST hiện đại. Các tài nguyên liên quan gồm
      \texttt{Patient}, \texttt{ServiceRequest}, \texttt{ImagingStudy} (tham
      chiếu tới ca chụp qua StudyInstanceUID và điểm truy cập DICOMweb) và
      \texttt{DiagnosticReport} \cite{hl7fhir}. FHIR phù hợp cho các hệ thống
      mới và cho trao đổi liên viện.
\item[Hồ sơ IHE] IHE không định nghĩa chuẩn mới mà quy định cách phối hợp các
      chuẩn trên trong từng kịch bản. Liên quan trực tiếp là Scheduled
      Workflow (SWF.b) cho luồng chỉ định--chụp--đọc, Patient Information
      Reconciliation (PIR) cho việc sửa định danh sau khi đã chụp, XDS-I.b cho
      chia sẻ ảnh giữa các cơ sở \cite{iherad}, và Audit Trail and Node
      Authentication (ATNA) cho xác thực nút mạng bằng TLS và định dạng bản ghi
      nhật ký truy vết \cite{iheiti}.
\end{description}

Mức độ tích hợp khả thi phụ thuộc bối cảnh. Ở bối cảnh A, HIS/RIS thường đã hỗ
trợ HL7 và worklist, và nền tảng chỉ cần nhận bản tin. Ở bối cảnh C, thường không
có RIS, worklist không khả dụng và định danh được nhập tay trên máy chụp; nền tảng
phải chấp nhận thực tế này và cung cấp chức năng đối chiếu, sửa định danh có ghi
nhật ký theo tinh thần PIR. Ngoài ra, Thông tư 13/2025/TT-BYT yêu cầu hồ sơ bệnh
án điện tử kết nối với cơ sở dữ liệu quốc gia về dân cư (Mục~\ref{sec:legal}); khi
đó số định danh cá nhân trở thành khoá liên kết bệnh nhân giữa các cơ sở, đồng thời
là một định danh trực tiếp phải bị loại bỏ khi khử định danh.

\subsection{Khử định danh}
\label{sec:deid}

Vì các định danh được nhúng ngay trong bản thân đối tượng ảnh, việc chia sẻ một
đối tượng DICOM đồng nghĩa với chia sẻ danh tính bệnh nhân, trừ khi có hành động
chủ động ngăn chặn. DICOM PS3.15 Phụ lục E \cite{dicomps315} định nghĩa các Hồ sơ bảo mật
(Confidentiality Profile) quy định thuộc tính nào phải bị loại bỏ, thay thế hoặc
giữ lại tùy theo mục đích sử dụng.

Ba khó khăn thường tái diễn trong thực tế và phải được xử lý tường minh:

\begin{enumerate}
    \item \textbf{Chú thích in chìm vào ảnh.} Một số thiết bị vẽ thông tin bệnh
          nhân trực tiếp vào mảng điểm ảnh. Thuộc tính
          \texttt{(0028,0301) BurnedInAnnotation} báo hiệu điều này nhưng không
          được điền nhất quán, nên có thể cần kiểm tra ở mức điểm ảnh.
    \item \textbf{Thẻ riêng của nhà sản xuất.} Các hãng lưu dữ liệu độc quyền
          trong những nhóm thẻ riêng với ngữ nghĩa không được công bố. Chính sách
          an toàn là loại bỏ toàn bộ thẻ riêng, chấp nhận mất các tham số thu
          nhận ảnh mà chúng mang theo.
    \item \textbf{Tính nhất quán của UID.} Thay thế UID một cách ngẫu nhiên sẽ
          phá vỡ quan hệ ca chụp--chuỗi ảnh--đối tượng ảnh. Khử định danh phải áp
          dụng một phép ánh xạ đơn ánh nhất quán để giữ nguyên cấu trúc phân cấp
          trong khi các định danh gốc không thể khôi phục được.
\end{enumerate}

Hơn nữa, khử định danh cần được hiểu là giảm thiểu rủi ro chứ không phải loại bỏ
rủi ro. Ảnh CT sọ não chứa cấu trúc khuôn mặt của bệnh nhân, và việc dựng bề mặt
từ một khối ảnh đã khử định danh có thể tạo ra một khuôn mặt nhận diện được; khả
năng này đã được ghi nhận trong y văn \cite{schwarz2019identification} và có nghĩa là dữ liệu hình ảnh phải được
đối xử như dữ liệu cá nhân ngay cả sau khi đã loại bỏ các thẻ. Giá trị pháp lý
của kết luận này theo pháp luật Việt Nam được bàn luận ở Mục~\ref{sec:legal}.
